\section{Bölüm sonu çözümlü problemler}

\begin{enumerate}
  \item \textbf{ANOVA tablosunu tamamlama.} $SS_A=240$, $SS_E=600$, $k=4$ ve $N=44$ olsun. Serbestlik derecelerini, ortalama kareleri ve $F$ değerini bulun.

  \textbf{Çözüm:} $sd_A=k-1=3$, $sd_E=N-k=40$'tır. $MS_A=240/3=80$, $MS_E=600/40=15$ ve $F=80/15=5{,}33$ olur. Toplam serbestlik derecesi 43'tür.

  \item \textbf{Eta-kare ve omega-kare.} Önceki problemde $SS_T=840$ olsun. $\eta^2$ ve $\omega^2$ değerlerini bulun.

  \textbf{Çözüm:} $\eta^2=SS_A/SS_T=240/840\approx0{,}286$'dır. $\omega^2=(SS_A-sd_A MS_E)/(SS_T+MS_E)=(240-3\cdot15)/(840+15)=195/855\approx0{,}228$ olur.

  \item \textbf{Takip yöntemi.} Grup büyüklükleri dengesiz ve standart sapmalar belirgin farklıyken genel test anlamlı bulunmuştur. Hangi genel ve ikili yöntemler uygundur?

  \textbf{Çözüm:} Varyans eşitliğine dayanmayan Welch ANOVA uygun genel testtir. Farkın yerini incelemek için Games--Howell ikili karşılaştırmaları kullanılabilir. Genel test tek başına bütün grupların birbirinden farklı olduğunu göstermez.
\end{enumerate}